Information in languages must be organized and presented in a manner that clearly and effectively conveys the speaker’s intention. The study of how information is structured in language is wide and encompasses syntax, semantics, pragmatics, and prosody (see Breen et al., 2012; Lambrecht, 1994; Roberts, 2012). One important area of information structure is focus or the information that is considered most important, relevant, new or contrastive (Kiss, 1998). Focus is believed to be a linguistic universal (Comrie, 1989). Yet, how focus is implemented varies across languages and can be constrained by a language’s phonology and morphosyntax (Kiss, 1998; Lambrecht, 1994). How speakers process focus in sentences remains a rich area of psycholinguistics research as it can reveal much about language and cognition. The methods in which focus has been researched range from behavioral tasks (e.g., Cutler & Fodor, 1979; Paterson et al., 1999) to event-related-potentials (e.g., Chen et al., 2014; Wang et al., 2011) to eye-tracking (e.g., Filik et al., 2005; Höhle et al., 2016).

Whereas web-based experimentation is transforming language research, the effects and limitations of the move from in-person testing to web-based testing are still being understood. To the best of our knowledge, using the internet to research focus in sentence processing, particularly with eye-tracking methods, has yet to be attempted. Recent work has documented how moving an eye-tracking experiment to a web-based platform can introduce various implications for study design (e.g., changes in instrument accuracy, measurement frame rate, instrument design, and data processing; see Bramlett & Wiener, 2024; Vos et al., 2022). In the current study, we test whether web-based eye-tracking can be used to (closely) replicate previously published findings, validate current web-based methods, and carry out exploratory studies on individual differences.

Focus in English only-sentences

In the present study, we “focus” on English preverbal only-sentences with varying positions of prosodic prominence. As an example, take the sentence, “Obama only vetoed the bill.” The scope of the focus particle ‘only’ can be associated with the verb ‘vetoed’ or the object ‘the bill.’ Semantic parsing, however, depends on which word(s) carries prosodic prominence, which is generally conveyed through an expanded F0 range, increased amplitude, and longer duration (i.e., nuclear pitch accent on the focal element(s); Breen et al., 2012; Gussenhoven, 1983).  If ‘the bill’ carries prosodic prominence, the listener will understand that Obama vetoed nothing else but the bill. In contrast, if ‘vetoed’ carries prosodic prominence, the listener will understand that Obama did nothing else to the bill other than veto it. 

Processing focus in only-sentences requires multiple levels of linguistic knowledge and serves as a valuable test case for understanding how first language (L1) and second language (L2) processing differ. Ge et al. (2021)--the study we closely replicate–examined how L1 and L2 English speakers process only-sentences in real time. The authors made use of the look and listen visual world paradigm in which a participant looks at images on screen while listening to spoken sentences. Importantly, the images on the screen represented the intended focus target or the alternative focus (i.e., a competitor). Ge et al. (2021) found that L1 Dutch-L2 English and L1 Cantonese-L2 English speakers showed patterns of eye movements that differed from those of L1 English speakers. In particular, L2 speakers showed delayed eye movements to the alternative of focus. The authors interpreted these differences as evidence for problematic integration of multiple interfaces (e.g., syntax-semantics, syntax-pragmatics) in real time. 

In the current study, we first set out to closely replicate Ge et al. (2021) using web-based eye-tracking and open materials and code. This serves as an important contribution to the field and tests the limits of web-based visual world paradigm research. We also extend Ge et al. (2021) by probing multiple individual differences to determine what, if any, behavioral measures serve as reliable predictors of L1 and L2 focus processing. 
